\documentclass[12pt]{article}
\usepackage[a4paper,margin=2.2cm]{geometry}
\usepackage{amsmath,amssymb,amsthm,mathtools}
\usepackage{hyperref}
\usepackage{cleveref}

% 中文（xelatex）
\usepackage{fontspec}
\usepackage{xeCJK}

\usepackage{fontspec}
\usepackage{xeCJK}
\setCJKmainfont{Noto Serif CJK SC}

\title{Linear Algebra}
\date{}
\begin{document}
\maketitle


\begin{flushleft}
    1. 2x2 矩阵 $\begin{bmatrix}
            a & b \\
            c & d
        \end{bmatrix}$ 把单位正方形变化平行四边形，两条边分别是向量$(a, c)$和$(b, d)$。
    用叉积公式推导这个平行四边形的有向面积, 证明它等于 ad-bc, 也就是行列式 \\
    解: \\
    因为题目需要使用叉积公式推导(建立于三维及以上) \\
    那么我们需要把 基向量 $\vec{v1} = [a, c]$ 和 $\vec{v2} = [b, d]$ 提维为: \\
    $\vec{v1} = [a, c, g]$ 和 $\vec{v2} = [b, d, z]$  其中 $g,z$ = 0 \\
    那么我们会得到如下： \\
    $\vec{v1} = [a, c, g]$ \\
    $\vec{v2} = [b, d, z]$ \\

    我们已知叉积会得到一个向量(垂直于v1,v2所构成的平面(也就是那个平行四边形由 0, v1, v2, v1 + v2 构成))\\
    所以我需要算出这个叉积向量并获取其模长,也就是这个平行四边形的面积 \\
    已知道叉积公式如下： \\
    $\vec{v1}  \times \vec{v2}$ =  $(cz - gd , gb - az, ad -cb)$ \\
    已知 $g,z = 0$
    所以公式可以简化: \\
    $\vec{v1}  \times \vec{v2}$ =  $(0 , 0, ad -cb)$ \\
    = $\vec{v1}  \times \vec{v2}$ =  $(0 , 0, ad- bc)$ \\
    所以我们的可以基于 $\vec{r} = [0,0,ad -bc]$ 算出模长: \\
    $\sqrt{(0)^{2} + (0)^{2} + (ad- bc)^{2}}$ \\
    = $|ad - bc|$ \\
    这里我们得到了一个绝对值,但我们丢失了方向 \\
    已知$\vec{r} = [0,0,ad -bc]$ 可以转换成 $\vec{r} = x*0 + y*0 + k(ad-bc)$ \\
    所以 如果 ad- bc 是否大于0 来确定方向 ad - bc  > 0 是 逆时针 反之是顺时针
\end{flushleft}

\subsection{特征向量}
\begin{flushleft}
    特征向量公式:
    $A\vec{u} - \lambda\vec{u} = \vec{0}$
    这个向量0 非常有意思
    向量0的形态是: $\begin{bmatrix}
            0 \\
            0 \\
            ..
        \end{bmatrix}$  \\
    在公式中这个向量0 as $v0_{1}$ 由两个向量减得到 在几何视角中两个向量相减是距离 例如 A-B 就是 B走到A的距离 \\
    也就是说这个向量0 就表达了 $A\vec{u}$ 和 $\lambda\vec{u}$ 是完全重合的 \\

    那么我们把 $A\vec{u} - \lambda\vec{u} = \vec{0}$ 代数变化下: \\
    = $(A - \lambda I)\vec{u} = \vec{0}$ 其中 I 是标准矩阵 $\begin{bmatrix}
            1 & 0 \\
            0 & 1
        \end{bmatrix}$ \\
    那么这个向量0 as $v0_{2}$   是什么意思呢? 我们把 $(A - \lambda I) = K$ 矩阵 也就是说 $ K \cdot \vec{u} = \vec{0}$ \\
    $\vec{u}$ 被压到了零点？ 这个 $\vec{0}$ 是 0点的意思 \\
    $v0_{1}$ 和 $v0_{2}$ 的表达性质是不一样的 \\
    那么我们现在来研究 K 是什么性质 ？ \\
    K 的表达式是:  \\
    $A - \lambda I$ 这是一个标准的向量减法, 它的结果是一个矩阵\\
    这个矩阵表示的是两个矩阵的差异 是一个差异矩阵 \\
    所以我们知道 K 是一个差异矩阵 \\
    那么一个 差异矩阵 * $\vec{u}$ = $\vec{0}$
    是否可以表述 $A\vec{u} - \lambda I\vec{u} = 0$ 是同一件事情呢? \\
    在代数式上 它很简单也很明显  $K\vec{u}$ =  $A\vec{u} - \lambda I\vec{u}$ \\
    是描述同一件事? \\
    答案是: 在这个例子中 代数和几何是描述同一个事。\\
    它们更像描述文件地址 一个是绝对坐标($v0_{1}$), 一个是相对坐标($v0_{2}$)\\

    我有一个关键的点 虽然能明白一些 但我感觉到很模糊 \\
    $A\vec{u} - \lambda I\vec{u} = \vec{0} $ -> 式子a \\
    问题1: 等于向量0代表是完全重叠 而非在一个方向上(它比在同一个方向要严苛的多)
    我们来看 矩阵$A$ 和 矩阵 $\lambda I$ \\
    如果让式子a成立 那么有下面明显的下面几种明显的情况 \\
    1: $A$ = $\lambda I$ = 0  \\
    2: $A$ = $\lambda I$ 或 $A$ - $\lambda I$  = 0  \\
    但从代数视角来说 det($A$ - $\lambda I$) = 0 并不满足呀 \\
    当然我们从几何性质来说 det($A$ - $\lambda I$) = 0 是降维性质 可能会把 $\vec{u}$ 压到0点但不一定会压到0点 只有 式子a成立才会被压到0点 \\

    我大概明白了一些 其实关键在一个很基础点上 我忽视了 我们要找的是两个未知两 一个是 $\lambda$ 一个  $\vec{u}$ \\
    同方向是$\vec{u}$和 $A\vec{u}$ \\
    但 $\lambda$ 是一个可拉伸收缩的控制值 当两个向量在同一个方向 那么其中一个向量 *$\lambda$ 就一定可以和另一个向量重叠 \\
\end{flushleft}

\paragraph{Task 1.1 投影矩阵:}
\begin{flushleft}
    $A = \begin{bmatrix}
            1 & 0 \\
            0 & 0 \\
        \end{bmatrix}$ \\
    1: 纯靠你在脑海中想象这个矩阵对二维空间的变换动作 \\
    一眼看去 我就会把 A 拆称两个基向量 $\vec{i} = \begin{bmatrix}
            1 \\
            0
        \end{bmatrix}$ 和 $\vec{j} = \begin{bmatrix}
            0 \\
            0
        \end{bmatrix}$ \\
    我们发现 $\vec{j}(y)$ 是 0向量 所以 我们很明白这是一个只会沿着x轴运动的\\
    它把所有点都拍到x轴上了 \\
    2: 代数式解: \\
    $(A - \lambda I) =  \begin{bmatrix}
            1 - \lambda & 0           \\
            0           & 0 - \lambda
        \end{bmatrix}$   \\
    = $det((A - \lambda I)) = 0$ \\
    = $(1 - \lambda)(-\lambda) - 0  = 0$ \\
    = $-\lambda + \lambda^{2} - 0 = 0$ \\
    = $\lambda^{2} - \lambda - 0 = 0 $
    = $\lambda^{2} = \lambda$ \\
    so: $\lambda = 1  \| \lambda = 0$ \\
    if: $\lambda = 1$: \\
    = $\begin{bmatrix}
            0 & 0  \\
            0 & -1 \\
        \end{bmatrix} @ \begin{bmatrix}
            x \\
            y
        \end{bmatrix}$ \\
    = \begin{align}
        0x + 0y = 0  \\
        0x + -1y = 0 \\
    \end{align}
    so: y = 0 ,x 可以是任意数 \\
    $\vec{v} = \begin{bmatrix}
            1 \\
            0
        \end{bmatrix}$
    if: $\lambda = 0$: \\
    = $\begin{bmatrix}
            1 & 0 \\
            0 & 0 \\
        \end{bmatrix} @ \begin{bmatrix}
            x \\
            y
        \end{bmatrix}$ \\
    = \begin{align}
        1x + 0y = 0 \\
        0x + 0y = 0 \\
    \end{align}
    so: x = 0, y 可以是任意数 \\
    $\vec{v} = \begin{bmatrix}
            0 \\
            1
        \end{bmatrix}$ \\
    我们找到了两个组
    \begin{align}
        \vec{u} = \begin{bmatrix}
                      1 \\
                      0
                  \end{bmatrix} \& \lambda = 1 \\
        \vec{u} = \begin{bmatrix}
                      0 \\
                      1
                  \end{bmatrix} \& \lambda = 0
    \end{align}
    运算后发现 全在x轴(原点也是x轴上的一点) 没有推翻我们的几何推断
\end{flushleft}

\paragraph{任务 1.2 镜像矩阵:}
\begin{flushleft}
    $B =  \begin{bmatrix}
            0 & 1 \\
            1 & 0 \\
        \end{bmatrix}$
    1: 纯靠你在脑海中想象这个矩阵对二维空间的变换动作 \\
    新的基i向量是0,1  新的基j向量是1,0 与原基向量相反(也就是对称关系)
    经典的沿着x = y 这条直线对称
    2: 代数式解: \\
    = $det(\begin{bmatrix}
            -\lambda & 1
            1        & -\lambda
        \end{bmatrix}) = 0$ \\
    = $(-\lambda)(-\lambda) - 1 = 0$ \\
    = $\lambda^{2} - 1 = 0$ \\
    = $\lambda^{2} = 1$  \\
    = $\sqrt{\lambda} = 1 $ \\
    so : $\lambda = 1 \| \lambda = -1$ \\
    if : $\lambda = 1 $:
    = $\begin{bmatrix}
            -1 , 1
            1  , -1
        \end{bmatrix}\begin{bmatrix}
            x \\
            y
        \end{bmatrix} = 0$ \\
    \begin{align}
        -x + y = 0 \\
        x + -y = 0
    \end{align} \\
    x,y相等 \\
    if : $\lambda = -1 $:
    = $\begin{bmatrix}
            1 , 1
            1  , 1
        \end{bmatrix}\begin{bmatrix}
            x \\
            y
        \end{bmatrix} = 0$ \\
    \begin{align}
        x + y = 0 \\
        x + y = 0
    \end{align} \\
    x,y为相反数 \\
    so: \\
    \begin{align}
        \vec{u} =  \begin{bmatrix}
                       1 \\
                       1
                   \end{bmatrix} \& \lambda = 1 \\
        \vec{u} =  \begin{bmatrix}
                       1 \\
                       -1
                   \end{bmatrix} \& \lambda = -1
    \end{align}
\end{flushleft}

\paragraph{Task 2.1}
\begin{flushleft}
    C = $\begin{bmatrix}
            3 & 0 \\
            0 & 3
        \end{bmatrix}$ \\
    1: 纯靠你在脑海中想象这个矩阵对二维空间的变换动作 \\
    部数增大(3)

    2: 算 $\lambda $ 是多少\\
    = $\begin{bmatrix}
            3 - \lambda & 0          \\
            0           & 3- \lambda
        \end{bmatrix}$ \\
    = $(3 - \lambda)(3- \lambda) = 0$ \\
    那么 $\lambda$ 只有一个解 就是3 \\

    3: 当把$\lambda $代入回 $C -\lambda I$,会得到什么矩阵？\\
    = $\begin{bmatrix}
            0 & 0 \\
            0 & 0
        \end{bmatrix}$ \\
    会得到0矩阵

    4：这说明在这个空间里，特征向量到底有几根 ？\\
    我们由2知道 这个矩阵是个0矩阵 \\
    因为 $\begin{bmatrix}
            0 & 0 \\
            0 & 0
        \end{bmatrix}\begin{bmatrix}
            x \\
            y
        \end{bmatrix} = \vec{0}$ 成立, x,y可以为 $\forall$ (也就是所有向量都是3的特征向量)\\
    根只有为3的一根
\end{flushleft}


\paragraph{Task 2.2}
\begin{flushleft}
    D = $\begin{bmatrix}
            1 & 1 \\
            0 & 1
        \end{bmatrix}$
    1: 纯靠你在脑海中想象这个矩阵对二维空间的变换动作 \\
    我们观察到j基量由 0,1 变成 1,1. 我们从几何的习惯理解 当x越大整体的趋势会越向右,所以这很典型是一个往右推的矩阵变化(注意应该没有上提
    因为y并没有增大)

    2: 解出它的特征值 $\lambda $ \\
    = $det(\begin{bmatrix}
            1 -\lambda & 1           \\
            0          & 1 - \lambda
        \end{bmatrix}) = 0$ \\
    \begin{align}
        (1- \lambda)x + y = 0 \\
        (1 - \lambda)y  = 0
    \end{align} \\
    我们解上面的方程组 $  (1 - \lambda)y  = 0 $  \\
    = $(y = 0 \& \forall \lambda )  \| (\lambda = 1 \& \forall y )$ \\
    if $(y = 0 \& \forall \lambda )$: \\
    那么 x = 0 \\
    $\vec{v} = \begin{bmatrix}
            0 \\
            0
        \end{bmatrix}$ -> 无意义 \\
    if $(\lambda = 1 \& \forall y )$: \\
    带入式子1 $(1- \lambda)x + y = 0$ \\
    我们可以确定 y = 0 , $\forall x$ \\
    $\vec{v} = \begin{bmatrix}
            1 \\
            0
        \end{bmatrix}$

    3: 你发现了什么异常？相比于普通的 2×2
    矩阵，这个矩阵在“寻找方向不变的向量”这件事上，遭遇了什么“先天残疾”？
    这个先天的残疾 我并不能理解 难道只有一个特征值解(1) ? 或者 特征向量 y 固定为0 ？
    还有一个问题：
    还是这个式子:  \\
    $A\vec{v} =\lambda \vec{v}$
    对于这个矩阵A 如果这个代数式子成立
    那么证明我们找到了特征向量和特征值 那么其他向量呢
    它们没有意义吗？
    这里我有个点无法理解
    我们知道 $\vec{v} = \begin{bmatrix}
            3 \\
            1
        \end{bmatrix}$ 是标准坐标系的表达
    而 $c1v1 + c2v2 = \vec{v} = \begin{bmatrix}
            3 \\
            1
        \end{bmatrix}$
    从这个代数式来看 它说在找也给在新的基底 存在[c1,c2] 等于 [3,1]
    这点我无法理解
    那么正确的式子应该是这样呀
    c1v1 + c2v2 = 3ei + 1ej 呀 同一个位置
\end{flushleft}

\begin{flushleft}
    我追求干净和简洁和逻辑顺畅的理解 我先说下我的理解 :
    我们先向量:
    $\vec{v} = \begin{bmatrix}
            3 \\
            1
        \end{bmatrix}$ \\
    这向量基于标准基的向量 它是简化的 完整代数表达形式是: \\
    $\vec{v} = 3 \cdot \begin{bmatrix}
            1 \\
            0
        \end{bmatrix} + 1 \cdot \begin{bmatrix}
            0 \\
            1
        \end{bmatrix}$ \\
    由上面的代数式我们很清楚的明白这里有一个累加的逻辑 先走3个 $\vec{b_i}$ + 1 $\vec{b_j}$ \\
    到达了一个目的 这个目的是 $\vec{v}$ 这个向量v 首先它是个地点 它在标准基下环境名称是 [3, 1] \\
    就像变量名和指针地址的关系 在另一个基环境下 这个变量名是可以变化的 但指针地址并不会改变/ \\

    那么变基操作呢? \\
    我们把 3, 1换成 c1,c2, 两个基向量换成v1, v2呢 \\
    那么代数式(表达同一个地方): \\
    c1v1 + c2v2  =  $ 3 \cdot \begin{bmatrix}
            1 \\
            0
        \end{bmatrix} + 1 \cdot \begin{bmatrix}
            0 \\
            1
        \end{bmatrix}$ \\
    我们现在遇到一个矩阵 A  它会改变基向量 这里有个很重要的逻辑细节: \\
    先换基底而非先矩阵运算 \\
    那么式子就会变成这样: \\
    $A(c1v1+ c2v2) = A(\begin{bmatrix}
            3 \\
            1
        \end{bmatrix}) $ \\
    那么现在我们已知道有两组未知变化c,v 正常情况下是无解的\\
    如果 A存在特征向量/值 \\
    那么式子可以变化 : \\
    $c1(Av1) + c2(Av2) = A(\begin{bmatrix}
            3 \\
            1
        \end{bmatrix}) $ \\
    = $c1(\lambda v1) + c2(\lambda_{2} v2) = A(\begin{bmatrix}
            3 \\
            1
        \end{bmatrix})$ \\
    那么这个代数式子表达什么？ \\
    如果一个矩阵存在特征向量/值 那么通过换基(基换成特征向量)操作 走[c1,c2]
    是可以到达同一个地方($A(\begin{bmatrix}
            3 \\
            1
        \end{bmatrix})$)的？
    但有个问题 这是一个代数式 也就是 c1 c2 必须存在解才能证明这个可能？

    这是我的一点 如果有任何不正确和模糊的地方 你需要指出来

    你刚才说到了一个非常有数学美感的一句话:
    一个n维矩阵 如果它不是残缺矩阵(存在n个特征向量 但它们同一个方向)
    那么它就一定能证明 $c1(\lambda v1) + c2(\lambda_{2} v2) = A(\begin{bmatrix}
            3 \\
            1
        \end{bmatrix})$ 是有解的 例如: 2维 就会撑起一个平面 只要你能撑起这个平面 那么你就一定可以走到指定的目的地(条条大路通罗马？)
\end{flushleft}

\paragraph{Task 3}
\begin{flushleft}
    一个$2 x 2$的矩阵M $\begin{bmatrix}
            a & b \\
            c & d \\
        \end{bmatrix}$: \\
    1: 当它作用在向量 $\vec{v1} = \begin{bmatrix}
            1 \\
            1
        \end{bmatrix}$时, 拉长2倍,方向不变.\\
    2: 当它作用到向量 $\vec{v2}  = \begin{bmatrix}
            -1 \\
            1
        \end{bmatrix}$时 压成零向量. \\

    构建一个方程组: \\
    \begin{align}
        M @ v1 = 2(v1)           \\
        M @ v2 = 0               \\
        a + b = 2 (x)            \\
        c + d = 2 (y)            \\
        -a + b  = 0              \\
        -c + d  = 0              \\
        a + b - 0 =  2 - (a - b) \\
        c + d - 0 = 2 - (c - d)  \\
        a + b  = 2 - a + b       \\
        c + d  = 2 - c + d       \\
        a + b + a - b = 2        \\
        c + d + c - d  = 2       \\
        2a = 2                   \\
        2c = 2                   \\
        a, c = 1                 \\
        b, d = 1                 \\
    \end{align}
    so :  M = $\begin{bmatrix}
            1 & 1 \\
            1 & 1 \\
        \end{bmatrix}$\\
    思考：如果M运算连续作用100次($M^{100}$) 空间会变成什么样？ 如果以v1,v2作为新坐标系
    事情会变的简单？ \\
    连续作用100次 空间会变成什么样? \\
    空间会以指数级拉大 因为 a, b, c, d 大于等于1
    如果以v1,v2作为新坐标系
    事情会变的简单? \\
    对于v1来说 这个扩大的系数稳定为2\\
    那么就可以简化为 $v1 * 2^{100}$ \\
    v2就更简单 直接是0 ？

\end{flushleft}

\paragraph{Task 4}
\begin{flushleft}
    R = $\begin{bmatrix}
            0 & -1 \\
            1 & 0
        \end{bmatrix}$ \\
    1: 解 $det(R- \Lambda I) = 0$. \\
    = $\begin{bmatrix}
            -\lambda & -1       \\
            1        & -\lambda
        \end{bmatrix}$  \\
    = $(-\lambda)^{2} + 1 = 0$ \\
    = $ \lambda ^{2} + 1 = 0 $ ? 实数无解 ？ \\
\end{flushleft}

回到了最基础的式子也是最重要的式子
$A\vec{v1} = \lambda \vec{v1}$

$A\vec{v1}$ 这个就是在原坐标系中的部分组成的概念！\\
$A\vec{v1} = \lambda \vec{v1}$ 而这个代数式就直接说明了
特征向量和其他向量的关系
是我笨了嘛 我虽然好像懂了 但依然感觉很模糊


\paragraph{Power-iteration}
\begin{flushleft}
    A = $\begin{bmatrix}
            a & b \\
            c & d
        \end{bmatrix}$ \\
    普通向量: $\vec{v} =\begin{bmatrix}
            x \\
            y
        \end{bmatrix}$ \\
    特征向量: $\vec{ev} = \begin{bmatrix}
            e \\
            k
        \end{bmatrix}$

    1: 我们需要证明出为什么在连续同样矩阵运算之后 其他向量会趋同于主特征向量? \\


\end{flushleft}

一个目的地 你可以 以多个规则到达 例如: 曲线, 翻转或 直线 .... \\
对于一个矩阵 它在标准xy坐标系中的变化规则 可能是曲线 \\
但如果我们把基底一换(非xy坐标) 那么它的变化规则 可能是直线 (如果是 这组基底就是 特征向量 $\lambda$ 只是幅度) \\
向量 $v = \begin{bmatrix}
        x \\
        y
    \end{bmatrix}$ \\
在标准系中: \\
A @ $\vec{v}$ \\
= $x\begin{bmatrix}
        1 \\
        0
    \end{bmatrix} + y\begin{bmatrix}
        0 \\
        1
    \end{bmatrix}$ \\
= $\vec{v_n}$ \\
这个 $\vec{v_n}$ 是个绝对坐标 不存在它在什么基底 \\

A 存在一组特征向量 e1, e2 我们对A执行换底: \\
$\vec{v}$  = \\
A@$\begin{bmatrix}
        \lambda 1 \\
        \lambda 2
    \end{bmatrix}$ \\
= $\lambda1 \cdot e1  + \lambda2 \cdot e2$  (如果一个平面存在两个向量 那么这个两个向量足以撑去这个平面)\\
那么上面这个表达式就可以表达一个非常重要的思想:
一个向量在变基的情况下是由两个基向量组织而成 ->分量被不均匀拉伸(分量是特征向量)

其实在标准环境下的$\vec{v}$ 和特征向量组环境下的 $\vec{v}$ 是同一个绝对地址？



\end{document}